% !TeX root = ../Tex/main.tex

\section{Introducation} \label{sub:methodology}
This second part of the paper is dedicated to test the analysis conducted by Professor Atrhur Spirling \parencite{spirlingDemocratizationLinguisticComplexity2016}. Actually, it provides good reasons to not stop doing research on the topic. In fact, it does not offer any decisive or comprehensive proof of any hypothesis. It suggest that slighlty moving the focus of regressions and inferece closer to the data of 1868 makes us change perpective on linguistic simplification.

%\vspace*{1em}
%\textit{Knowledge of Institutional Setting:} The analysis conducted is not supported by knowledge of institutional setting. Under, this aspect it well reflects what Spirling stated \parencite{spirlingDemocratizationLinguisticComplexity2016} and does not go deeper. A broader review of the literature on hystorical context in UK and UK parliament is set to be part of future researches that will expand this research proposal.

\vspace*{1em}
\textit{Knowledge of Language:} A limit of the present proposal is linguistic knowledge. In order to strenghthen the results of causal inference it is good to have deep understanding of the language used at the time. This would also shape the selection of n-grams, discriding irrelevant term and valuing relevant expressions. %\textcolor{red}{ma forse vale la pena aggiungere qualche reference sulla lingua}.

\vspace*{1em}
\textit{Data:} When I first tried to reproduce \textcite{spirlingDemocratizationLinguisticComplexity2016}'s replication package I found out that it is incomplete. I texed Professor \textcite{spirlingDemocratizationLinguisticComplexity2016} and he confirmed that the first part of the pipeline is missing.

The first aim of the analysis is to produce that first part of the pipeline and test whether the results lead to the same conclusion. I already retrieved data from the \href{https://hansard.parliament.uk/}{UK Parliament website} and I have already computed the basic linguistic measures.

\vspace*{1em}
\textit{DGP and Theoretical Model:} I intend to lay down a game theory model to describe and formalise the structure of the debate inside the Parliament. I will follow the most relevant examples in the literature such as \textcite{prokschPoliticsParliamentaryDebate2014, backPoliticalPartiesParliaments2016}.
The theoretical model will provide the basis for the Data Generatig Process.

\vspace*{1em}
\textit{Causal Inference:} In this section I present the intuition that challenges the results of \textcite{spirlingDemocratizationLinguisticComplexity2016}. According to Professor Spirling analysis, after the broadening of electorate base of UK's House of Commons the language used in parliamentary speeches suffered of simplification in terms of lenghth of sentences and words used. In particular, \textcite{spirlingDemocratizationLinguisticComplexity2016} tests the hypothesis that the treatment affects cabinet members language more than non-cabinet members language.

In the first part of the causal inference section I will challenge the hypothesis that the treatment affects at all the language in the parliament. I will use a DinD to estimate the coefficient of the the difference between linguistic trend in the two housesof the UK Parliament. Since the house of Lords is a non-elective chamber, the widening of the electorate should not influence the linguistic measures of the debate in this chamber. Thus, the measures of this house should provide a good control group (if parallel trend assumption holds).

In the second part I will build a topic model to create variables that describe to what extent a debate in parliament belongs to a specif topic. I will use these variables to understand how different topics influence the dependent variable.

In the third part I narrow down the window of data and I challenge \textcite{spirlingDemocratizationLinguisticComplexity2016}'s main hypothesis, namely the treatment affects cabinet members language more than non-cabinet members language.
I still have to find support for this test. I will deep my knowledge on bandwidth selection to understand the main drivers of this choice and its main effects on the resulting coefficients.

In the fourth part I will use topic model to understand whether the hypothesis to topic specification.

While in the original papaer from \textcite{spirlingDemocratizationLinguisticComplexity2016}, they use only one linguistic measure. I intend to provide the results from more statistics, to avoid biases that may arise from a single measure.
 

\section{Causal Inference}
%\lipsum[1] % genera il primo paragrafo di testo Lorem Ipsum

\subsection{RDD}
The identification strategy emplyed in this proposal is regression discontinuity design. Among the classical literature on RDD, I like highlighting two great contributions: \textcite{fujiwaraVotingTechnologyPolitical2015, gulzarPoliticiansBureaucratsDevelopment2017}. A more recent study on RDD methodology is presented by \textcite{cattaneoPracticalIntroductionRegression2024}.

The equation used is the following:

\vspace*{0.3em}
\[
FRE = \beta_0 + \beta_1 \cdot Cutoff + \gamma X + \varepsilon
\]
\vspace*{0.1em}
\[
FRE = \beta_0 + \beta_1 \cdot C + \gamma X + \varepsilon
\]

\vspace*{1em}
Identification strategy relies on:
\begin{itemize}
    \item[-] \small{\(C\)} is the treatment. It is an exgoneous variable independet from the others called cutoff. In our regression the cutoff is 1868, namely the year in which the reform was introduced
    \item[-] \small{\(Y\)} is any of the measures of linguistic simplification employed.
\end{itemize}

\vspace*{1em}
\textit{Controls:} The vector of controls is \small{\(X\)}, while \small{\(\gamma\)} is the vector of coefficients. It represents the effects of control variables on our outcome variable. The selected controls are:
\begin{itemize}
    \item[-] Competitiveness (of electoral distric)
    \item[-] Party
\end{itemize}
These variables are the same as \textcite{spirlingDemocratizationLinguisticComplexity2016}.


%------- End of the File -------